% =========================================================
% Functions: Definitions and Basic Properties
% =========================================================

\subsection{Functions}

% ---------------------------------------------------------
% TOOLKIT
% ---------------------------------------------------------

\begin{remark*}[Exposition]
A function is a deterministic machine: every input in the domain is accepted,
and no input can fork into two different outputs. Fibers are the level sets
sitting over points of the codomain; they collect all inputs that the machine
sends to the same output.
\end{remark*}

% ---------------------------------------------------------
% Function definition
% ---------------------------------------------------------
\begin{definitionbox}{Definition (Functional Relation)}
\begin{definition}[Functional Relation]\label{def:functional-relation}
Let $R\subseteq A\times B$ be a relation. We say that $R$ is
\emph{functional from $A$ to $B$} if each input has at most one output:
\[
\forall a\in A\,\forall b_1,b_2\in B,\quad
\bigl((a,b_1)\in R\land(a,b_2)\in R\bigr)\Rightarrow b_1=b_2.
\]
\end{definition}
\end{definitionbox}

\begin{remark*}[Standard quantified statement]
\[
  \operatorname{FunctionalRelation}(R,A,B)
  \iff
  \forall a\in A\,\forall b_1,b_2\in B,\ 
  ((a,b_1)\in R\wedge(a,b_2)\in R)\to b_1=b_2.
\]
\end{remark*}

\begin{remark*}[Predicate reading]
\[
  \operatorname{FunctionalRelation}(R,A,B)
  \iff
  \forall a\in A\,\forall b_1,b_2\in B,\ 
  ((a,b_1)\in R\wedge(a,b_2)\in R)\to b_1=b_2.
\]
\end{remark*}

\begin{remark*}[Interpretation]
A functional relation assigns at most one output to each input.
\end{remark*}

\begin{dependencies}
\begin{itemize}
  \item \hyperref[def:relation]{Relation}
  \item \hyperref[def:cartesian-product-rel]{Cartesian Product}
\end{itemize}
\end{dependencies}

\begin{definitionbox}{Definition (Total Relation on a Domain)}
\begin{definition}[Total Relation on a Domain]\label{def:total-relation-on-domain}
Let $R\subseteq A\times B$ be a relation. We say that $R$ is \emph{total on
$A$} if each input has at least one output:
\[
\forall a\in A\,\exists b\in B,\quad (a,b)\in R.
\]
\end{definition}
\end{definitionbox}

\begin{remark*}[Standard quantified statement]
\[
  \operatorname{TotalOnDomain}(R,A,B)
  \iff
  \forall a\in A,\ \exists b\in B,\ (a,b)\in R.
\]
\end{remark*}

\begin{remark*}[Predicate reading]
\[
  \operatorname{TotalOnDomain}(R,A,B)
  \iff
  \forall a\in A,\ \exists b\in B,\ (a,b)\in R.
\]
\end{remark*}

\begin{remark*}[Interpretation]
Totality on a domain means every declared input has at least one output.
\end{remark*}

\begin{dependencies}
\begin{itemize}
  \item \hyperref[def:relation]{Relation}
\end{itemize}
\end{dependencies}

\begin{definitionbox}{Definition (Function)}
\begin{definition}[Function]\label{def:function}
Let $A$ and $B$ be sets. A \emph{function} from $A$ to $B$ is a relation
$f \subseteq A \times B$, together with the declared domain $A$ and codomain
$B$, such that:
\begin{enumerate}[label=(\roman*)]
\item (\emph{Existence / left-total}) for every $a \in A$, there exists $b \in B$
  with $(a,b) \in f$;
\item (\emph{Uniqueness / right-unique}) if $(a,b_1) \in f$ and $(a,b_2) \in f$
  then $b_1 = b_2$.
\end{enumerate}
If $(a,b) \in f$ we write $f(a) = b$.
\end{definition}
\end{definitionbox}

\begin{remark*}[Standard quantified statement]
\[
  \operatorname{Function}(f,A,B)
  \iff
  \forall a\in A\,\exists! b\in B,\ (a,b)\in f.
\]
\end{remark*}

\begin{remark*}[Predicate reading]
\[
  \operatorname{Function}(f,A,B)
  \iff
  \operatorname{TotalOnDomain}(f,A,B)\wedge
  \operatorname{FunctionalRelation}(f,A,B).
\]
\end{remark*}

\begin{remark*}[Interpretation]
A function is a total and single-valued relation from its domain to its
codomain.
\end{remark*}

\begin{remark*}[Exposition]
A function is a rule assigning to each input exactly one output. The two
conditions formalize ``defined everywhere'' (existence) and ``single-valued''
(uniqueness).
\end{remark*}

\begin{remark*}[Exposition]
Every function is a relation, but not every relation is a function. The two
conditions that distinguish functions are left-totality and right-uniqueness.
The codomain is part of the function data, not something recoverable from the
graph alone.
\end{remark*}

\begin{dependencies}
\begin{itemize}
  \item \hyperref[def:relation]{Relation}
  \item \hyperref[def:cartesian-product-rel]{Cartesian Product}
\end{itemize}
\end{dependencies}

\begin{theorembox}{Theorem (Functional Relation Characterization of Functions)}
\begin{theorem}[Functional Relation Characterization of Functions]\label{thm:function-total-functional-characterization}
Let $A$ and $B$ be sets, and let $R\subseteq A\times B$. Then $R$ is a
function from $A$ to $B$ if and only if
\[
\forall a\in A\,\exists! b\in B,\quad (a,b)\in R.
\]
Equivalently, $R$ is a function from $A$ to $B$ if and only if $R$ is total
on $A$ and functional from $A$ to $B$.
\hyperref[prf:function-total-functional-characterization]{Go to proof.}
\end{theorem}
\end{theorembox}

\begin{remark*}[Standard quantified statement]
\[
  R\text{ is a function }A\to B
  \iff
  \forall a\in A\,\exists! b\in B,\ (a,b)\in R.
\]
\end{remark*}

\begin{remark*}[Predicate reading]
\[
  \operatorname{Function}(R,A,B)
  \iff
  \operatorname{TotalOnDomain}(R,A,B)\wedge
  \operatorname{FunctionalRelation}(R,A,B).
\]
\end{remark*}

\begin{remark*}[Interpretation]
The total-and-functional relation conditions are exactly the usual
unique-output condition for functions.
\end{remark*}

\begin{dependencies}
\begin{itemize}
  \item \hyperref[def:function]{Function}
  \item \hyperref[def:functional-relation]{Functional Relation}
  \item \hyperref[def:total-relation-on-domain]{Total Relation on a Domain}
\end{itemize}
\end{dependencies}

\begin{definitionbox}{Definition (Partial Function)}
\begin{definition}[Partial Function]\label{def:partial-function}
Let $A$ and $B$ be sets. A \emph{partial function} from $A$ to $B$ is a
functional relation $R\subseteq A\times B$. It need not be total on $A$:
some inputs may have no output.
\end{definition}
\end{definitionbox}

\begin{remark*}[Standard quantified statement]
\[
  \operatorname{PartialFunction}(R,A,B)
  \iff
  R\subseteq A\times B\wedge \operatorname{FunctionalRelation}(R,A,B).
\]
\end{remark*}

\begin{remark*}[Interpretation]
A partial function is allowed to be undefined at some inputs but remains
single-valued where it is defined.
\end{remark*}

\begin{dependencies}
\begin{itemize}
  \item \hyperref[def:functional-relation]{Functional Relation}
\end{itemize}
\end{dependencies}

\begin{propositionbox}{Proposition (Partial Functions Are Functional Relations)}
\begin{proposition}[Partial Functions Are Functional Relations]\label{prop:partial-function-functional-relation}
For sets $A$ and $B$ and a relation $R\subseteq A\times B$, $R$ is a partial
function from $A$ to $B$ if and only if $R$ is functional from $A$ to $B$.
\hyperref[prf:partial-function-functional-relation]{Go to proof.}
\end{proposition}
\end{propositionbox}

\begin{remark*}[Standard quantified statement]
\[
  \operatorname{PartialFunction}(R,A,B)
  \iff
  \operatorname{FunctionalRelation}(R,A,B).
\]
\end{remark*}

\begin{remark*}[Interpretation]
Partial functions are exactly functional relations when totality is not
required.
\end{remark*}

\begin{dependencies}
\begin{itemize}
  \item \hyperref[def:partial-function]{Partial Function}
  \item \hyperref[def:functional-relation]{Functional Relation}
\end{itemize}
\end{dependencies}

\begin{definitionbox}{Definition (Function Equality)}
\begin{definition}[Function Equality]\label{def:function-equality}
For functions $f:A\to B$ and $g:C\to D$, we write $f=g$ if
\[
A=C,\qquad B=D,\qquad\text{and}\qquad
\forall a\in A,\; f(a)=g(a).
\]
\end{definition}
\end{definitionbox}

\begin{remark*}[Standard quantified statement]
\[
  f=g
  \iff
  A=C\wedge B=D\wedge \forall a\in A,\ f(a)=g(a).
\]
\end{remark*}

\begin{remark*}[Predicate reading]
\[
  \operatorname{FunctionEqual}(f,g)
  \iff
  \operatorname{DomainOfFunction}(f,A)\wedge
  \operatorname{DomainOfFunction}(g,A)\wedge
  \operatorname{CodomainOfFunction}(f,B)\wedge
  \operatorname{CodomainOfFunction}(g,B)\wedge
  \forall a\in A,\ f(a)=g(a).
\]
\end{remark*}

\begin{remark*}[Interpretation]
Function equality requires agreement of domain, codomain, and all values.
\end{remark*}

\begin{dependencies}
\begin{itemize}
  \item \hyperref[def:function]{Function}
  \item \hyperref[def:set-equality]{Set Equality}
\end{itemize}
\end{dependencies}

\begin{propositionbox}{Proposition (One-to-Many Relations Are Not Functional)}
\begin{proposition}[One-to-Many Relations Are Not Functional]\label{prop:one-to-many-not-functional}
If $R\subseteq A\times B$ is one-to-many, then $R$ is not functional from $A$
to $B$. Consequently, $R$ is not a function from $A$ to $B$.
\hyperref[prf:one-to-many-not-functional]{Go to proof.}
\end{proposition}
\end{propositionbox}

\begin{remark*}[Standard quantified statement]
\[
  \operatorname{OneToManyRelation}(R,A,B)\to
  \neg\operatorname{FunctionalRelation}(R,A,B).
\]
\end{remark*}

\begin{remark*}[Interpretation]
One input with two distinct outputs violates the uniqueness condition for
functions.
\end{remark*}

\begin{remark*}[Exposition]
A function may send different inputs to the same output. That is the ordinary
many-one case in Tarski's terminology. What is forbidden for a function is not
shared output, but one input having two distinct outputs.
\end{remark*}

\begin{dependencies}
\begin{itemize}
  \item \hyperref[def:one-to-many-relation]{One-to-Many Relation}
  \item \hyperref[def:functional-relation]{Functional Relation}
  \item \hyperref[def:function]{Function}
\end{itemize}
\end{dependencies}

\begin{definitionbox}{Definition (Domain of a Function)}
\begin{definition}[Domain of a Function]\label{def:domain}
If $f$ is a function from $A$ to $B$, we write $f : A \to B$, where $A$ is
the \emph{domain} $\dom(f)$.
\end{definition}
\end{definitionbox}

\begin{remark*}[Standard quantified statement]
\[
  \operatorname{DomainOfFunction}(f,A)
  \iff
  f:A\to B\text{ for some }B.
\]
\end{remark*}

\begin{remark*}[Interpretation]
The domain is the declared set of admissible inputs.
\end{remark*}

\begin{dependencies}
\begin{itemize}
  \item \hyperref[def:function]{Function}
\end{itemize}
\end{dependencies}

\begin{definitionbox}{Definition (Codomain of a Function)}
\begin{definition}[Codomain of a Function]\label{def:codomain}
If $f$ is a function from $A$ to $B$, we write $f : A \to B$, where $B$ is
the \emph{codomain} $\cod(f)$.
\end{definition}
\end{definitionbox}

\begin{remark*}[Standard quantified statement]
\[
  \operatorname{CodomainOfFunction}(f,B)
  \iff
  f:A\to B\text{ for some }A.
\]
\end{remark*}

\begin{remark*}[Interpretation]
The codomain is the declared target set, whether or not every target is hit.
\end{remark*}

\begin{remark*}[Exposition]
The notation $f:A\to B$ records both the domain $A$ and the codomain $B$.
The same graph may define maps into larger codomains, but surjectivity and
bijectivity depend on the declared codomain.
\end{remark*}

\begin{dependencies}
\begin{itemize}
  \item \hyperref[def:function]{Function}
\end{itemize}
\end{dependencies}

\begin{definitionbox}{Definition (Image of a Function)}
\begin{definition}[Image of a Function]\label{def:image-function}
The \emph{image} (or \emph{range}) of $f : A \to B$ is:
\[
\im(f) \;:=\; \{\, b \in B \mid \exists a \in A,\ f(a)=b \,\}.
\]
The image may be a proper subset of the codomain.
\end{definition}
\end{definitionbox}

\begin{remark*}[Standard quantified statement]
\[
  b\in\im(f)
  \iff
  \exists a\in A,\ f(a)=b.
\]
\end{remark*}

\begin{remark*}[Predicate reading]
\[
  \operatorname{ImageElement}(b,f,A)
  \iff
  \exists a\in A,\ f(a)=b.
\]
\end{remark*}

\begin{remark*}[Interpretation]
The image is the subset of the codomain actually attained by the function.
\end{remark*}

\begin{remark*}[Exposition]
For a function viewed as its graph, the relation range and the function image
coincide. The codomain may be larger than this range/image.
\end{remark*}

\begin{remark*}[Exposition]
The chapter \hyperref[ch:functions-and-order]{Functions and Order} uses the
image as a replacement codomain: every function is surjective onto its image,
and an injective function becomes bijective after corestricting to its image.
\end{remark*}

\begin{dependencies}
\begin{itemize}
  \item \hyperref[def:function]{Function}
  \item \hyperref[def:set-membership]{Set and Membership}
  \item \hyperref[ax:separation]{Axiom Schema of Separation}
\end{itemize}
\end{dependencies}

\begin{definitionbox}{Definition (Image of a Set)}
\begin{definition}[Image of a Set]\label{def:image-set}
For $S \subseteq A$, the \emph{image of $S$ under $f$} is:
\[
f(S) \;:=\; \{\, f(a) \mid a \in S \,\}.
\]
Note $f(A) = \im(f)$.
\end{definition}
\end{definitionbox}

\begin{remark*}[Standard quantified statement]
\[
  y\in f(S)
  \iff
  \exists a\in S,\ f(a)=y.
\]
\end{remark*}

\begin{remark*}[Interpretation]
The image of a subset collects the values attained on that subset.
\end{remark*}

\begin{dependencies}
\begin{itemize}
  \item \hyperref[def:function]{Function}
  \item \hyperref[def:subset]{Subset}
  \item \hyperref[ax:separation]{Axiom Schema of Separation}
\end{itemize}
\end{dependencies}

\begin{definitionbox}{Definition (Preimage)}
\begin{definition}[Preimage]\label{def:preimage}
For $T \subseteq B$, the \emph{preimage} (inverse image) of $T$ under $f$ is:
\[
f^{-1}(T) \;:=\; \{\, a \in A \mid f(a) \in T \,\}.
\]
\end{definition}
\end{definitionbox}

\begin{remark*}[Standard quantified statement]
\[
  a\in f^{-1}(T)
  \iff
  a\in A\wedge f(a)\in T.
\]
\end{remark*}

\begin{remark*}[Predicate reading]
\[
  \operatorname{PreimageElement}(a,f,T)
  \iff
  f(a)\in T.
\]
\end{remark*}

\begin{remark*}[Interpretation]
The preimage pulls a set of target values back to the inputs that land in it.
\end{remark*}

\begin{remark*}[Exposition]
$f^{-1}(T)$ does \emph{not} denote an inverse function. It is defined for any
function $f$, regardless of whether $f$ is injective or bijective.
It sends a subset of the codomain back to a subset of the domain. An inverse
function, when it exists, sends individual elements of the codomain back to
individual elements of the domain.
\end{remark*}

\begin{dependencies}
\begin{itemize}
  \item \hyperref[def:function]{Function}
  \item \hyperref[def:subset]{Subset}
  \item \hyperref[ax:separation]{Axiom Schema of Separation}
\end{itemize}
\end{dependencies}

\begin{figure}[htbp]
\centering
\begin{tikzpicture}[atlas, x=1cm, y=1cm, >=Stealth, every node/.style={font=\small}]
  \node[atlas label] at (-2.1,2.65) {Preimage of a set};
  \draw[atlasgray, line width=0.7pt] (-3.85,-0.1) ellipse (0.75 and 1.65);
  \draw[atlasgray, line width=0.7pt] (-1.55,-0.1) ellipse (0.7 and 1.65);
  \draw[atlasgold, line width=1pt, rounded corners=8pt] (-2.05,0.05) rectangle (-1.05,1.0);
  \node[atlas label] at (-3.85,1.75) {$A$};
  \node[atlas label] at (-1.55,1.75) {$B$};
  \node[atlas label] at (-1.05,1.15) {$T$};

  \node[atlas dot, label={[atlas label,left]$a_1$}] (pa1) at (-3.85,0.85) {};
  \node[atlas dot, label={[atlas label,left]$a_2$}] (pa2) at (-3.85,0.05) {};
  \node[atlas dot, label={[atlas label,left]$a_3$}] (pa3) at (-3.85,-1.05) {};
  \node[atlas dot, label={[atlas label,right]$b_1$}] (pb1) at (-1.55,0.55) {};
  \node[atlas dot, label={[atlas label,right]$b_2$}] (pb2) at (-1.55,-1.05) {};

  \draw[atlasblue, atlas probe, ->] (pa1) -- (pb1);
  \draw[atlasblue, atlas probe, ->] (pa2) -- (pb1);
  \draw[atlasblue, atlas probe, ->] (pa3) -- (pb2);
  \draw[atlasgreen, line width=1pt, rounded corners=8pt] (-4.42,-0.17) rectangle (-3.28,1.08);
  \node[atlas label] at (-3.85,-1.95) {$f^{-1}(T)\subseteq A$};

  \node[atlas label] at (2.0,2.65) {Inverse function};
  \draw[atlasgray, line width=0.7pt] (0.25,-0.1) ellipse (0.7 and 1.45);
  \draw[atlasgray, line width=0.7pt] (2.45,-0.1) ellipse (0.7 and 1.45);
  \node[atlas label] at (0.25,1.55) {$A$};
  \node[atlas label] at (2.45,1.55) {$B$};

  \node[atlas dot, label={[atlas label,left]$a_1$}] (ia1) at (0.25,0.65) {};
  \node[atlas dot, label={[atlas label,left]$a_2$}] (ia2) at (0.25,-0.85) {};
  \node[atlas dot, label={[atlas label,right]$b_1$}] (ib1) at (2.45,0.65) {};
  \node[atlas dot, label={[atlas label,right]$b_2$}] (ib2) at (2.45,-0.85) {};

  \draw[atlasblue, atlas probe, ->] (ia1) -- (ib1);
  \draw[atlasblue, atlas probe, ->] (ia2) -- (ib2);
  \draw[atlasred, atlas probe, ->, bend left=18] (ib1) to (ia1);
  \draw[atlasred, atlas probe, ->, bend right=18] (ib2) to (ia2);
  \node[atlas label] at (1.35,-1.95) {$f^{-1}:B\to A$};
\end{tikzpicture}

\caption{A preimage of a set is not the same object as an inverse function.}
\label{fig:preimage-vs-inverse}
\end{figure}

\begin{definitionbox}{Definition (Fiber)}
\begin{definition}[Fiber]\label{def:fiber}
The \emph{fiber} of $f$ over $b \in B$ is the preimage of the singleton:
\[
f^{-1}(\{b\}) \;=\; \{\, a \in A \mid f(a) = b \,\}.
\]
\end{definition}
\end{definitionbox}

\begin{remark*}[Standard quantified statement]
\[
  a\in f^{-1}(\{b\})
  \iff
  f(a)=b.
\]
\end{remark*}

\begin{remark*}[Interpretation]
A fiber is the set of all inputs with one fixed output value.
\end{remark*}

\begin{remark*}[Exposition]
The collection $\{f^{-1}(\{b\}) \mid b \in \im(f)\}$ is a partition of $A$.
Every function thus induces an equivalence relation $a_1 \sim_f a_2 \iff f(a_1) = f(a_2)$,
whose equivalence classes are precisely the fibers.
\end{remark*}

\begin{dependencies}
\begin{itemize}
  \item \hyperref[def:function]{Function}
  \item \hyperref[def:preimage]{Preimage}
\end{itemize}
\end{dependencies}

\begin{definitionbox}{Definition (Graph of a Function)}
\begin{definition}[Graph of a Function]\label{def:graph}
The \emph{graph} of $f : A \to B$ is:
\[
\operatorname{Graph}(f)
\;:=\;
\{\, (a,b) \in A \times B \mid b = f(a) \,\}.
\]
\end{definition}
\end{definitionbox}

\begin{remark*}[Standard quantified statement]
\[
  (a,b)\in\operatorname{Graph}(f)
  \iff
  a\in A\wedge b=f(a).
\]
\end{remark*}

\begin{remark*}[Interpretation]
The graph records every input together with its unique output.
\end{remark*}

\begin{remark*}[Exposition]
A relation $G \subseteq A \times B$ is the graph of a function $f : A \to B$
iff $G$ is left-total and right-unique. Functions may therefore be identified
with their graphs only after the domain and codomain have been fixed. As bare
sets of ordered pairs, graphs do not encode the codomain.
\end{remark*}

\begin{dependencies}
\begin{itemize}
  \item \hyperref[def:function]{Function}
  \item \hyperref[def:ordered-pair]{Ordered Pair}
  \item \hyperref[def:cartesian-product]{Cartesian Product}
\end{itemize}
\end{dependencies}

% ---------------------------------------------------------
% Injectivity / Surjectivity / Bijectivity
% ---------------------------------------------------------
\begin{definitionbox}{Definition (Injective Function)}
\begin{definition}[Injective Function]\label{def:injective}
$f : A \to B$ is \emph{injective} (one-to-one) if distinct elements of $A$
have distinct images:
\[
\forall a_1,a_2 \in A,\;
f(a_1)=f(a_2) \;\Longrightarrow\; a_1=a_2.
\]
\end{definition}
\end{definitionbox}

\begin{remark*}[Standard quantified statement]
\[
  \operatorname{Injective}(f)
  \iff
  \forall a_1,a_2\in A,\ (f(a_1)=f(a_2)\to a_1=a_2).
\]
\end{remark*}

\begin{remark*}[Predicate reading]
\[
  \operatorname{Injective}(f)
  \iff
  \forall a_1,a_2\in A,\ (f(a_1)=f(a_2)\to a_1=a_2).
\]
\end{remark*}

\begin{remark*}[Interpretation]
Injectivity means that equal outputs can only come from equal inputs.
\end{remark*}

\begin{dependencies}
\begin{itemize}
  \item \hyperref[def:function]{Function}
\end{itemize}
\end{dependencies}

\begin{definitionbox}{Definition (Empty Function)}
\begin{definition}[Empty Function]\label{def:empty-function}
For every set $B$, the \emph{empty function} from $\varnothing$ to $B$ is the
empty relation
\[
\varnothing:\varnothing\to B.
\]
It is the unique function with domain $\varnothing$ and codomain $B$.
\end{definition}
\end{definitionbox}

\begin{remark*}[Standard quantified statement]
\[
  \operatorname{EmptyFunction}(\varnothing,B).
\]
\end{remark*}

\begin{remark*}[Interpretation]
The empty function has no inputs and therefore no value assignments.
\end{remark*}

\begin{dependencies}
\begin{itemize}
  \item \hyperref[def:function]{Function}
  \item \hyperref[def:empty-set]{Empty Set}
  \item \hyperref[def:codomain]{Codomain of a Function}
\end{itemize}
\end{dependencies}

\begin{definitionbox}{Definition (Surjective Function)}
\begin{definition}[Surjective Function]\label{def:surjective}
$f : A \to B$ is \emph{surjective} (onto) if every element of $B$ is achieved:
\[
\forall b \in B,\; \exists a \in A \text{ such that } f(a)=b.
\]
Equivalently, $\im(f) = B$.
\end{definition}
\end{definitionbox}

\begin{remark*}[Standard quantified statement]
\[
  \operatorname{Surjective}(f)
  \iff
  \forall b\in B,\ \exists a\in A,\ f(a)=b.
\]
\end{remark*}

\begin{remark*}[Predicate reading]
\[
  \operatorname{Surjective}(f)
  \iff
  \forall b\in B,\ \exists a\in A,\ f(a)=b.
\]
\end{remark*}

\begin{remark*}[Interpretation]
Surjectivity means every declared target is hit by at least one input.
\end{remark*}

\begin{dependencies}
\begin{itemize}
  \item \hyperref[def:function]{Function}
\end{itemize}
\end{dependencies}

\begin{definitionbox}{Definition (Bijective Function)}
\begin{definition}[Bijective Function]\label{def:bijective}
$f : A \to B$ is \emph{bijective} if it is both injective and surjective.
\end{definition}
\end{definitionbox}

\begin{remark*}[Standard quantified statement]
\[
  \operatorname{Bijective}(f)
  \iff
  \operatorname{Injective}(f)\wedge\operatorname{Surjective}(f).
\]
\end{remark*}

\begin{remark*}[Interpretation]
Bijectivity means that each target has exactly one input mapping to it.
\end{remark*}

\begin{remark*}[Exposition]
Injectivity: each fiber has at most one element. Surjectivity: every fiber is
nonempty. Bijectivity: every fiber has exactly one element.
\end{remark*}

\begin{dependencies}
\begin{itemize}
  \item \hyperref[def:function]{Function}
  \item \hyperref[def:injective]{Injective Function}
  \item \hyperref[def:surjective]{Surjective Function}
\end{itemize}
\end{dependencies}

\begin{propositionbox}{Proposition (One-One Relations and Converse Functionality)}
\begin{proposition}[One-One Relations and Converse Functionality]\label{prop:one-to-one-converse-functional}
Let $f:A\to B$ be a function, identified with its graph
$G_f\subseteq A\times B$. Then $f$ is injective if and only if the converse
relation $G_f^{-1}\subseteq B\times A$ is functional from $B$ to $A$.
\hyperref[prf:one-to-one-converse-functional]{Go to proof.}
\end{proposition}
\end{propositionbox}

\begin{remark*}[Standard quantified statement]
\[
  \operatorname{Injective}(f)
  \iff
  \operatorname{FunctionalRelation}(G_f^{-1},B,A).
\]
\end{remark*}

\begin{remark*}[Interpretation]
Injectivity of a function is exactly functionality of its converse graph.
\end{remark*}

\begin{dependencies}
\begin{itemize}
  \item \hyperref[def:injective]{Injective Function}
  \item \hyperref[def:converse-relation]{Converse Relation}
  \item \hyperref[def:functional-relation]{Functional Relation}
  \item \hyperref[def:graph]{Graph of a Function}
\end{itemize}
\end{dependencies}

\begin{propositionbox}{Proposition (Bijections Have Unique Correspondence Both Ways)}
\begin{proposition}[Bijections Have Unique Correspondence Both Ways]\label{prop:bijection-unique-both-directions}
Let $f:A\to B$ be a function. Then $f$ is bijective if and only if
\[
\forall a\in A\,\exists! b\in B,\quad f(a)=b,
\]
and
\[
\forall b\in B\,\exists! a\in A,\quad f(a)=b.
\]
\hyperref[prf:bijection-unique-both-directions]{Go to proof.}
\end{proposition}
\end{propositionbox}

\begin{remark*}[Standard quantified statement]
\[
  \operatorname{Bijective}(f)
  \iff
  \forall b\in B\,\exists!a\in A,\ f(a)=b.
\]
\end{remark*}

\begin{remark*}[Predicate reading]
\[
  \operatorname{Bijective}(f)
  \iff
  \operatorname{Injective}(f)\wedge\operatorname{Surjective}(f).
\]
\end{remark*}

\begin{remark*}[Interpretation]
A bijection gives unique correspondence from inputs to outputs and from outputs
back to inputs.
\end{remark*}

\begin{dependencies}
\begin{itemize}
  \item \hyperref[def:bijective]{Bijective Function}
  \item \hyperref[def:injective]{Injective Function}
  \item \hyperref[def:surjective]{Surjective Function}
\end{itemize}
\end{dependencies}

\begin{definitionbox}{Definition (Many-Place Functional Relation)}
\begin{definition}[Many-Place Functional Relation]\label{def:many-place-functional-relation}
A relation $R\subseteq A_1\times\cdots\times A_n\times B$ is
\emph{functional in the last coordinate} if
\[
\forall a_1\in A_1\cdots\forall a_n\in A_n\,
\forall b_1,b_2\in B,
\]
\[
\bigl((a_1,\ldots,a_n,b_1)\in R\land
(a_1,\ldots,a_n,b_2)\in R\bigr)\Rightarrow b_1=b_2.
\]
When such a relation is also total in the first $n$ coordinates, it determines
a function of $n$ variables, written $f(a_1,\ldots,a_n)=b$.
\end{definition}
\end{definitionbox}

\begin{remark*}[Standard quantified statement]
\[
  \forall a_1\in A_1\cdots\forall a_n\in A_n\,\forall b_1,b_2\in B,\quad
  ((a_1,\ldots,a_n,b_1)\in R\wedge(a_1,\ldots,a_n,b_2)\in R)\to b_1=b_2.
\]
\end{remark*}

\begin{remark*}[Predicate reading]
\[
  \operatorname{ManyPlaceFunctionalRelation}(R,A_1,\ldots,A_n,B)
  \iff
  \forall \vec a\,\forall b_1,b_2\in B,\
  ((\vec a,b_1)\in R\wedge(\vec a,b_2)\in R)\to b_1=b_2.
\]
\end{remark*}

\begin{remark*}[Interpretation]
A many-place functional relation is single-valued in its output coordinate.
\end{remark*}

\begin{dependencies}
\begin{itemize}
  \item \hyperref[def:many-place-relation]{Many-Place Relation}
  \item \hyperref[def:functional-relation]{Functional Relation}
\end{itemize}
\end{dependencies}

% ---------------------------------------------------------
% Special functions
% ---------------------------------------------------------
\begin{definitionbox}{Definition (Identity Function)}
\begin{definition}[Identity Function]\label{def:identity}
The \emph{identity function} on $A$ is $\id_A : A \to A$, $\id_A(a) = a$.
\end{definition}
\end{definitionbox}

\begin{remark*}[Standard quantified statement]
\[
  \forall a\in A,\quad \id_A(a)=a.
\]
\end{remark*}

\begin{remark*}[Interpretation]
The identity function returns each input unchanged.
\end{remark*}

\begin{dependencies}
\begin{itemize}
  \item \hyperref[def:function]{Function}
\end{itemize}
\end{dependencies}

\begin{definitionbox}{Definition (Inclusion Map)}
\begin{definition}[Inclusion Map]\label{def:inclusion}
If $A \subseteq B$, the \emph{inclusion map} is $\iota : A \hookrightarrow B$,
$\iota(a) = a$.
\end{definition}
\end{definitionbox}

\begin{remark*}[Standard quantified statement]
\[
  A\subseteq B\to \forall a\in A,\ \iota(a)=a.
\]
\end{remark*}

\begin{remark*}[Interpretation]
The inclusion map views elements of a subset as elements of the larger set.
\end{remark*}

\begin{remark*}[Exposition]
Both send each element to itself, but the inclusion map has a strictly larger
codomain when $A \subsetneq B$. The inclusion map is always injective.
\end{remark*}

\begin{dependencies}
\begin{itemize}
  \item \hyperref[def:function]{Function}
  \item \hyperref[def:subset]{Subset}
\end{itemize}
\end{dependencies}

\begin{definitionbox}{Definition (Constant Function)}
\begin{definition}[Constant Function]\label{def:constant}
$f : A \to B$ is \emph{constant} if there exists $b_0 \in B$ with
$f(a) = b_0$ for all $a \in A$.
\end{definition}
\end{definitionbox}

\begin{remark*}[Standard quantified statement]
$f$ is constant $\;\leftrightarrow\; \exists b_0 \in B,\; \forall a \in A,\; f(a) = b_0$.
\end{remark*}

\begin{remark*}[Predicate reading]
\[
  \operatorname{ConstantFunction}(f)
  \iff
  \exists b_0\in B,\ \forall a\in A,\ f(a)=b_0.
\]
\end{remark*}

\begin{remark*}[Interpretation]
A constant function sends every input to one fixed output.
\end{remark*}

\begin{dependencies}
\begin{itemize}
  \item \hyperref[def:function]{Function}
\end{itemize}
\end{dependencies}
